% Summary of: Coelho, Franco, Firoozabadi (2025)
% "Phase Equilibria of CO2-Water and CO2-Brine at High Temperatures:
%  From Monte Carlo Simulations to the Equation of State"
% Ind. Eng. Chem. Res. 64(16), 8492-8505. DOI: 10.1021/acs.iecr.5c00134
%
% This file is a plain-text LaTeX reference summary for use during
% our own paper writing. It captures all key equations, parameters,
% and results from the paper. Open Access (CC-BY 4.0).

\documentclass[11pt]{article}
\usepackage{amsmath, amssymb}
\usepackage{booktabs}
\usepackage{hyperref}
\usepackage[margin=2.5cm]{geometry}

\title{\textbf{Summary: Coelho, Franco \& Firoozabadi (2025)}\\[4pt]
\large Phase Equilibria of CO$_2$--Water and CO$_2$--Brine at High Temperatures:\\
From Monte Carlo Simulations to the Equation of State\\[4pt]
\normalsize \textit{Ind.\ Eng.\ Chem.\ Res.}\ \textbf{64}(16), 8492--8505\\
DOI: \href{https://doi.org/10.1021/acs.iecr.5c00134}{10.1021/acs.iecr.5c00134}}
\author{Summary compiled for the SALTbasis/Claude\_code project}
\date{2026-03-16}

\begin{document}
\maketitle
\tableofcontents

% ==============================================================================
\section{Overview and Scope}

\paragraph{Goal.} Parametrise CPA and eCPA equations of state (EoS) to describe
\begin{itemize}
  \item CO$_2$--H$_2$O phase equilibria (VLE + LLE) from 283 to 623~K and 5 to 3500~bar,
  \item CO$_2$ solubility in NaCl brine from 283 to 723~K, 1--1400~bar, 0--6~mol\,kg$^{-1}$,
  \item aqueous and CO$_2$-rich phase densities,
  \item Gibbs ensemble Monte Carlo (GEMC) simulations at high $T$ (548--623~K) as independent validation.
\end{itemize}

\paragraph{Key advance over previous work.} Including CO$_2$--H$_2$O \emph{cross-association}
(Wertheim theory) and \emph{temperature-dependent} binary interaction parameters reduces the
average CO$_2$ solubility deviation in brine from 26.6\% (Maribo-Mogensen et al.\ 2015)
to 6.9\% --- a 74\% improvement.

% ==============================================================================
\section{eCPA Equation of State}

The residual Helmholtz free energy is a sum of four contributions:
\begin{equation}
  A^r = A^{\text{phys}} + A^{\text{assoc}} + A^{\text{DH}} + A^{\text{Born}}
  \label{eq:Ar}
\end{equation}
When no electrolytes are present, $A^{\text{DH}} = A^{\text{Born}} = 0$ and the CPA EoS
is recovered.

\subsection{Physical term (SRK cubic)}
\begin{equation}
  \frac{A^{\text{SRK}}}{NRT} = -\ln\!\left(1 - \frac{b}{v}\right)
    - \frac{a}{bRT}\ln\!\left(1 + \frac{b}{v}\right)
\end{equation}

\paragraph{Mixing rules -- neutral components (van der Waals):}
\begin{align}
  b &= \sum_i b_i x_i \\
  a &= \sum_i \sum_j x_i x_j a_{ij}(1-k_{ij}), \qquad a_{ij} = \sqrt{a_i a_j}
\end{align}

\paragraph{Mixing rules -- electrolyte systems (Huron--Vidal infinite-pressure):}
\begin{equation}
  \frac{a}{b} = \sum_i x_i \frac{a_i}{b_i} - \frac{g^{E\infty}}{\ln 2}
\end{equation}
\begin{equation}
  g^{E\infty} = \frac{1}{b}\sum_j\sum_i x_i x_j b_j \Delta U_{ji}
\end{equation}
The $\Delta U_{ji}$ are NRTL-like energy parameters (zero non-randomness parameter).
For neutral--neutral pairs the two approaches are equivalent via
\begin{equation}
  \frac{\Delta U_{ji}}{\ln 2} = \frac{a_i}{b_i} - \frac{2\sqrt{a_i a_j}}{b_i+b_j}(1-k_{ij})
\end{equation}

\subsection{Association term (Wertheim)}
\begin{equation}
  \frac{A^{\text{assoc}}}{NRT}
    = \sum_i x_i \sum_{A\in i}
      \left(\ln\chi^{A_i} - \frac{\chi^{A_i}}{2} - \frac{1}{2}\right)
\end{equation}
where $\chi^{A_i}$ is the fraction of molecules $i$ \emph{not} bonded at site $A$:
\begin{align}
  \chi^{A_i} &= \frac{1}{1 + c\sum_j x_j \sum_{B\in j} \chi^{B_j}\,\delta^{A_iB_j}} \\
  \Delta^{A_iB_j} &= g_\eta\,\kappa^{A_iB_j}\frac{P}{RT}
    \left[\exp\!\left(\frac{\varepsilon^{A_iB_j}}{k_BT}\right)-1\right]
    = \frac{P}{RT}\,\delta^{A_iB_j}
\end{align}
\begin{equation}
  g_\eta = \frac{1}{1 - 1.9\eta}, \qquad \eta = \frac{b}{4v}
\end{equation}

\paragraph{H$_2$O association scheme.} 4C (2 proton donors + 2 proton acceptors).

\paragraph{CO$_2$--H$_2$O cross-association.}
CO$_2$ is treated as a pseudo-associating species with the 4C scheme; its cross-association
strength with water is proportional to the water self-association:
\begin{equation}
  \delta^{A_wB_c} = S_{cw}\,\delta^{A_wB_w}
\end{equation}
The cross-association is symmetric ($\delta^{A_wB_c} = \delta^{B_wA_c}$), so the
non-bonded fractions reduce to:
\begin{align}
  \chi_w &= \frac{1}{1 + 2c\,x_w\chi_w\delta_w + 2c\,x_c\chi_c S_{cw}\delta_w}
  \label{eq:chi_w}\\
  \chi_c &= \frac{1}{1 + 2c\,x_w\chi_w S_{cw}\delta_w}
  \label{eq:chi_c}
\end{align}

\subsection{Debye--H\"uckel term}
\begin{equation}
  \frac{A^{\text{DH}}}{NRT}
    = -\sum_j x_j z_j^2 \frac{X_j}{4\pi N_A c\sum_j x_j z_j^2}
\end{equation}
\begin{equation}
  X_i = \frac{1}{\sigma_i^3}
    \left[\ln(1+\kappa\sigma_i) - \kappa\sigma_i + \frac{(\kappa\sigma_i)^2}{2}\right]
\end{equation}
\begin{equation}
  \kappa = \sqrt{\frac{e^2 c \sum_j x_j z_j^2}{k_B T\,\varepsilon_0\varepsilon_r}}
  \qquad \text{(inverse Debye length)}
\end{equation}

\subsection{Born term}
\begin{equation}
  \frac{A^{\text{Born}}}{NRT}
    = \frac{N_A e^2}{8\pi\varepsilon_0 RT}
      \left(\frac{1}{\varepsilon_r} - 1\right)
      \sum_j x_j z_j^2 R_{B,j}
\end{equation}
$R_{B,j}$ is the Born cavity radius (differs from the hard-sphere diameter $\sigma_j$
to account for solvation).

\subsection{Relative permittivity model (primitive, Onsager--Kirkwood)}
\begin{align}
  \frac{(2\varepsilon_r+\varepsilon_\infty)(\varepsilon_r-\varepsilon_\infty)}
       {\varepsilon_r(\varepsilon_\infty+2)^2}
    &= \frac{N_A}{9\varepsilon_0 k_B T v} \sum_j x_j g_j \mu_{0j}^2
  \label{eq:permittivity}\\
  \frac{\varepsilon_\infty-1}{\varepsilon_\infty+2}
    &= \frac{N_A}{3\varepsilon_0 v}\sum_j x_j \alpha_{0j}
  \quad \text{(Clausius--Mossotti)}
\end{align}
The Kirkwood $g$-factor links the dielectric constant to the Wertheim association:
\begin{equation}
  g_i = 1 + \sum_j z_{ij} P_{ij}
    \cos\gamma_{ij}\frac{P_i\cos\theta_{ij}+1}{\mu_{0j}/\mu_{0i}}
\end{equation}
\begin{equation}
  P_{ij} = N_{A_i} c\, x_j\, \chi^{A_i} \chi^{B_j}\, \delta^{A_iB_j}
  \qquad \text{(probability of association)}
\end{equation}

% ==============================================================================
\section{Parametrisation}

\subsection{Pure component parameters}
Taken from Tsivintzelis et al.\ (2011). Temperature dependence of the SRK energy parameter:
\begin{equation}
  a_i = a_{0i}\left[1 + c_{1i}\!\left(1 - \sqrt{\frac{T}{T_{c,i}}}\right)\right]^2
\end{equation}
Ion parameters from Maribo-Mogensen et al.\ (2015): $a_i^{\text{ion}}=0$,
$b_i = 2N_A\pi\sigma_i^3/3$. A Peneloux volume translation is applied for brine
density ($v = v^{\text{EoS}} + x_s\Upsilon_s$).

\subsection{CO$_2$--H$_2$O binary parameters (temperature-dependent)}
\begin{align}
  k_{cw} &= A_k\!\left(\frac{T}{T_{c,c}}\right)^2 + B_k\!\left(\frac{T}{T_{c,c}}\right) + C_k
  \label{eq:kcw}\\
  S_{cw} &= A_S\!\left(\frac{T}{T_{c,c}}\right)^2 + B_S\!\left(\frac{T}{T_{c,c}}\right) + C_S
  \label{eq:Scw}
\end{align}
$T_{c,c} = 304$~K (CO$_2$ critical temperature).

\begin{table}[h]
  \centering
  \caption{Optimised CPA parameters for CO$_2$--H$_2$O (Eqs.~\ref{eq:kcw}--\ref{eq:Scw})}
  \begin{tabular}{lrrr}
    \toprule
    Parameter & $A$ & $B$ & $C$ \\ \midrule
    $k_{cw}$  & $-0.492$ & $2.101$ & $-1.571$ \\
    $S_{cw}$  & $0.192$  & $-0.173$ & $-0.009$ \\
    \bottomrule
  \end{tabular}
\end{table}

\subsection{CO$_2$--NaCl interaction parameters}
The interaction between neutral species $i$ and salt $s$ is temperature-dependent:
\begin{equation}
  \Delta U_{is} = \Delta U_{is}^{\text{ref}}
    + \alpha_{is} R\!\left[
      \left(1 - \frac{T}{T_{i,s}^\alpha}\right)^2
      - \left(1 - \frac{T_{\text{ref}}}{T_{i,s}^\alpha}\right)^2
    \right], \quad T_{\text{ref}} = 298\,\text{K}
  \label{eq:DeltaU}
\end{equation}

\begin{table}[h]
  \centering
  \caption{Optimised eCPA interaction parameters for CO$_2$--NaCl--H$_2$O (Eq.~\ref{eq:DeltaU})}
  \begin{tabular}{lrrr}
    \toprule
    Parameter & $\Delta U_{is}^{\text{ref}}$ (J\,mol$^{-1}$) & $\alpha_{is}$ (K) & $T_{i,s}^\alpha$ (K) \\ \midrule
    $\Delta U_{ws}$ (water--NaCl)  & $-1858$ & $1573$ & $340$ \\
    $\Delta U_{cs}$ (CO$_2$--NaCl) & $+6056$ & $692$  & $244$ \\
    \bottomrule
  \end{tabular}
\end{table}

The water--NaCl parameter ($\Delta U_{ws}$) is taken from Maribo-Mogensen et al.\ (2015).
The CO$_2$--NaCl parameter ($\Delta U_{cs}$) is fitted to CO$_2$ solubility in brine.

\subsection{Objective functions}
\paragraph{CPA (CO$_2$--H$_2$O):}
\begin{equation}
  \text{O.F.} = \min \sum_k
    \left(\frac{1}{N_k^W}\sum_i\left|\frac{x_{c,i}^{W,\text{calc}}-x_{c,i}^{W,\text{exp}}}{x_{c,i}^{W,\text{exp}}}\right|
    + \frac{1}{N_k^C}\sum_i\left|\frac{x_{c,i}^{C,\text{calc}}-x_{c,i}^{C,\text{exp}}}{x_{c,i}^{C,\text{exp}}}\right|\right)
\end{equation}
Temperature subsets with upper limits: $\{298, 323, 353, 393, 423, 473, 538, 573, 598, 623\}$~K.

\paragraph{eCPA (CO$_2$--NaCl brine):}
\begin{equation}
  \text{O.F.} = \min \sum_i \left|\frac{m_{c,i}^{\text{calc}} - m_{c,i}^{\text{exp}}}{m_{c,i}^{\text{exp}}}\right|
\end{equation}

% ==============================================================================
\section{Performance Summary}

\subsection{CPA: CO$_2$--H$_2$O}
Fit range: 283--623~K, 5--3500~bar. Global AAD:
\begin{itemize}
  \item Aqueous phase ($x_c^W$): \textbf{8.6\%} over 539 data points
  \item CO$_2$-rich phase ($x_c^C$): \textbf{23.9\%} over 443 data points
\end{itemize}
Density of aqueous phase (293--468~K, 1--800~bar): \textbf{0.47\%} AAD. \\
Density of CO$_2$-rich phase (304--478~K, 1--1300~bar): \textbf{5.33\%} AAD.

\subsection{eCPA: CO$_2$--NaCl brine}
Fit range: 283--723~K, 1--1400~bar, 0.17--6~mol\,kg$^{-1}$. \\
Global AAD for CO$_2$ solubility: \textbf{6.9\%} over 660 data points.\\
Improvement over Maribo-Mogensen et al.\ (2015): 26.6\% $\to$ 6.9\% (\textbf{74\% reduction}).\\
Density of CO$_2$-saturated brine: \textbf{1.05\%} AAD.

\paragraph{Caution.} The model may not correctly represent the solubility trend at $T < 298$~K
(trend with temperature is incorrect in this range).

% ==============================================================================
\section{Component Numbering Convention (eCPA ELV)}

\textbf{Important for implementation.} In the eCPA ELV equilibrium equations and in
the \texttt{ecpa/} code, the component index convention is:

\begin{center}
\begin{tabular}{lll}
  \toprule
  Index & Component & Phase \\ \midrule
  1 ($x_{1w}$) & H$_2$O & aqueous \\
  2 ($x_{2w}$) & Na$^+$ & aqueous \\
  3 ($x_{3w}$) & Cl$^-$ & aqueous \\
  4 ($x_{4w}$) & CO$_2$ & aqueous \\
  1 ($x_{1c}$) & H$_2$O & CO$_2$-rich \\
  4 ($x_{4c} = 1-x_{1c}$) & CO$_2$ & CO$_2$-rich \\
  \bottomrule
\end{tabular}
\end{center}

This is the \emph{opposite} of the CPA2 (salt-free) labeling where component 1 is CO$_2$.
Key conversions for comparing with experimental data:
\begin{align}
  m_c \; [\text{mol\,kg}^{-1}] &= \frac{x_{4w}}{x_{1w} M_w} \\
  x_c^W \; [\text{salt-incl.}] &= x_{4w} \\
  x_c^W \; [\text{salt-free}] &= \frac{x_{4w}}{x_{4w} + x_{1w}} \\
  x_c^C &= x_{4c} = 1 - x_{1c}
\end{align}
where $M_w = 0.018015$~kg\,mol$^{-1}$.

% ==============================================================================
\section{Monte Carlo Simulations}

\paragraph{Setup.} NPT-GEMC, Cassandra software. 2000 molecules, two-box system.
50M steps equilibration + 50M steps production. CO$_2$ global mole fraction: 15--30\%.

\paragraph{Force fields.}
\begin{itemize}
  \item Water: SPCE (rigid)
  \item CO$_2$: EPM2 (semiflexible: rigid bonds, harmonic angle potential)
  \item Na$^+$, Cl$^-$: classical integer-charge FF
\end{itemize}
Unlike LJ parameters: Lorentz--Berthelot (LB) or optimised (OPT) from Vlcek et al.\ (2011).
LJ cutoff 1.2~nm with analytical tail corrections; Ewald summation ($10^{-5}$ accuracy).

\paragraph{Key finding.} SPCE-EPM2 with OPT unlike parameters reproduces the high-$T$
(548--623~K) CO$_2$--H$_2$O phase diagram without polarization corrections.
At high $T$, thermal effects overcome SPCE overpolarization.

% ==============================================================================
\section{Key References from the Paper}

\begin{description}
  \item[eCPA framework] Maribo-Mogensen et al.\ (2015) --- original eCPA with DH + Born terms
  \item[CPA pure params] Tsivintzelis et al.\ (2011) --- CO$_2$ and H$_2$O SRK/CPA parameters
  \item[Cross-association] Li \& Firoozabadi (2009) [RERI] --- CO$_2$--H$_2$O cross-association in CPA
  \item[Unlike LJ params] Vlcek et al.\ (2011) --- OPT SPCE-EPM2 parameters for GEMC
  \item[High-$T$ exp.\ CO$_2$--H$_2$O] Tödheide \& Franck (1963); Takenouchi \& Kennedy (1964)
  \item[High-$T$ exp.\ CO$_2$--brine] Takenouchi \& Kennedy (1965)
  \item[Low-$T$ exp.\ CO$_2$--brine] Guo et al.\ (2015); Rumpf et al.\ (1994); Kiepe et al.\ (2002); Koschel et al.\ (2006)
\end{description}

% ==============================================================================
\section{Conclusions (verbatim from paper)}

\begin{enumerate}
  \item New CPA parameters for CO$_2$--H$_2$O describe the phase diagram accurately from
    283--623~K and for $T > 304$~K in both phases.
  \item eCPA with CO$_2$--H$_2$O cross-association and temperature-dependent parameters
    reduces CO$_2$ solubility deviation in brine by 74\% vs.\ previous model.
    The ternary mixture uses four binary parameters (quadratic polynomial in $T$).
  \item Aqueous phase density deviation $\approx 1\%$; CO$_2$ dissolution may increase
    or decrease brine density depending on salinity and temperature.
  \item SPCE-EPM2 (OPT) is a good force-field combination for high-$T$ GEMC simulations.
  \item CO$_2$--H$_2$O cross-association confirmed from both EoS and molecular structure
    (radial distribution functions, coordination numbers).
\end{enumerate}

% ==============================================================================
\section{Supporting Information}

\subsection*{Table S1: Pure Component Parameters}

\begin{center}
\begin{tabular}{lcccccccccccc}
\hline
 & $a_0$ & $c_1$ & $b$ & $\varepsilon/k_B$ & $\kappa$ & $\sigma$ & $R_B$ & $\mu_0$ & $\alpha_0$ & $\gamma$ & $\theta$ \\
 & (J\,m³/mol²) & & (cm³/mol) & (K) & & (Å) & (Å) & (D) & (ų) & (mN/m) & (deg) \\
\hline
H$_2$O & 0.1228 & 0.6736 & 14.52 & 2003 & 1.004 & 1.855 & 2.36 & 1.67 & 1.613 & 63.5 & 94.8 \\
CO$_2$ & 0.3508 & 0.7602 & 27.20 & --- & --- & --- & --- & --- & --- & --- & --- \\
Na$^+$ & --- & --- & 16.49 & --- & --- & 3.19 & 1.83 & --- & --- & --- & --- \\
Cl$^-$ & --- & --- & 40.83 & --- & --- & 3.557 & --- & --- & --- & --- & --- \\
\hline
\end{tabular}
\end{center}

\noindent NaCl Peneloux volume shift: $\Upsilon = -53.5$~cm$^3$/mol.

Parameters:
\begin{itemize}
  \item $a_0$, $c_1$: SRK energy parameter and Mathias-Copeman constant
  \item $b$: co-volume (cm³/mol)
  \item $\varepsilon/k_B$: association energy / Boltzmann constant (K); H$_2$O uses 4C scheme
  \item $\kappa$: association volume; H$_2$O only
  \item $\sigma$: hard-sphere diameter for DH term (Å)
  \item $R_B$: Born radius (Å)
  \item $\mu_0$: dipole moment (Debye); H$_2$O only
  \item $\alpha_0$, $\gamma$, $\theta$: polarizability (ų), surface tension (mN/m), H-O-H angle (deg); H$_2$O only
\end{itemize}

\subsection*{Table S2: MD Force-Field Parameters}

\paragraph{SPC/E water (rigid):}
\begin{center}
\begin{tabular}{lccc}
\hline
Site & $\sigma$ (Å) & $\varepsilon$ (kJ/mol) & $q$ (e) \\
\hline
O$_w$ & 3.166 & 0.6502 & $-0.8476$ \\
H$_w$ & 0 & 0 & $+0.4238$ \\
\hline
\end{tabular}
\end{center}
O--H bond: 1.0 Å; H--O--H angle: 109.47°.

\paragraph{EPM2 CO$_2$ (semiflexible: rigid bonds, harmonic angle):}
\begin{center}
\begin{tabular}{lccc}
\hline
Site & $\sigma$ (Å) & $\varepsilon$ (kJ/mol) & $q$ (e) \\
\hline
C$_c$ & 2.757 & 0.2339 & $+0.6512$ \\
O$_c$ & 3.033 & 0.6693 & $-0.3256$ \\
\hline
\end{tabular}
\end{center}
C--O bond: 1.149 Å; O--C--O angle equilibrium: 180°; angle force constant $k_\theta = 1236$~kJ\,mol$^{-1}$\,rad$^{-2}$.

\paragraph{Na$^+$ and Cl$^-$:} Classical integer charges ($q = \pm1$).
From Joung \& Cheatham (2008) adapted for SPC/E.

\subsection*{Table S3: Unlike Lennard-Jones Parameters}

\begin{center}
\begin{tabular}{lccccc}
\hline
Pair & \multicolumn{2}{c}{Lorentz--Berthelot (LB)} & \multicolumn{2}{c}{Optimised (OPT)} \\
 & $\sigma$ (Å) & $\varepsilon$ (kJ/mol) & $\sigma$ (Å) & $\varepsilon$ (kJ/mol) \\
\hline
C$_c$--O$_w$ & 2.962 & 0.389 & 2.840 & 0.550 \\
O$_c$--O$_w$ & 3.100 & 0.660 & 3.162 & 0.660 \\
\hline
\end{tabular}
\end{center}

\noindent OPT parameters from Vlcek et al.\ (2011). The OPT parameters reduce the C$_c$--O$_w$ distance
by $\sim$4\% and increase the well depth by $\sim$41\% relative to LB, improving CO$_2$ solubility
predictions at high temperature.

\subsection*{Supporting Equations (Fugacity Coefficients)}

The compressibility factor for a mixture is:
\begin{equation}
  Z = \frac{PV}{NkT} = 1 + \frac{B(\mathbf{x},T)}{V} + \cdots
\end{equation}

The residual Helmholtz energy per mole is:
\begin{equation}
  A^\mathrm{res} = A^\mathrm{phys} + A^\mathrm{assoc} + A^\mathrm{DH} + A^\mathrm{Born}
\end{equation}

The fugacity coefficient of component $i$ in the mixture:
\begin{equation}
  \ln \hat{\phi}_i = \frac{1}{RT}\left(\frac{\partial (nA^\mathrm{res})}{\partial n_i}\right)_{T,V,n_{j\neq i}}
    - \ln Z
\end{equation}

\paragraph{Physical term (SRK):}
\begin{align}
  \frac{A^\mathrm{phys}}{NkT} &= -\ln\!\left(1 - \frac{b}{V}\right)
    - \frac{a}{bRT}\ln\!\left(1 + \frac{b}{V}\right) \\
  \ln\hat{\phi}_i^\mathrm{phys} &= \frac{b_i}{b}(Z-1) - \ln\!\left(Z - \frac{Pb}{RT}\right)
    - \frac{a}{bRT}\left(\frac{2\sum_j x_j a_{ij}}{a} - \frac{b_i}{b}\right)
    \ln\!\left(1 + \frac{Pb}{ZRT}\right)
\end{align}

\paragraph{Association term:}
\begin{equation}
  \frac{A^\mathrm{assoc}}{NkT} = \sum_i x_i \sum_{A_i}
    \left[\ln X^{A_i} - \frac{X^{A_i}}{2} + \frac{1}{2}\right]
\end{equation}
where $X^{A_i}$ is the fraction of site $A$ on molecule $i$ not bonded:
\begin{equation}
  X^{A_i} = \frac{1}{1 + \rho \sum_j x_j \sum_{B_j} X^{B_j} \Delta^{A_i B_j}}
\end{equation}
\begin{equation}
  \Delta^{A_i B_j} = g^\mathrm{hs}(\sigma_{ij}) \left[\exp\!\left(\frac{\varepsilon^{A_i B_j}}{k_B T}\right) - 1\right]
    \kappa^{A_i B_j} b_{ij}
\end{equation}

\paragraph{Debye--Hückel term:}
\begin{equation}
  \frac{A^\mathrm{DH}}{NkT} = -\frac{\kappa^3 V}{3\pi N_A}
  \cdot \frac{1}{1 + \kappa \bar{a}}
  \cdot \frac{1}{4\pi\varepsilon_0 \varepsilon_r k_B T}
\end{equation}
with inverse screening length:
\begin{equation}
  \kappa^2 = \frac{e^2}{\varepsilon_0 \varepsilon_r k_B T} \sum_i \rho_i z_i^2
\end{equation}
The DH contribution to $\ln\hat{\phi}_i$ involves derivatives of $\kappa$ and $\varepsilon_r$ w.r.t.\ composition.

\paragraph{Born term:}
\begin{equation}
  \frac{A^\mathrm{Born}}{NkT} = -\frac{e^2}{8\pi\varepsilon_0 k_B T}
    \sum_i \frac{n_i z_i^2}{R_{B,i}}\left(\frac{1}{\varepsilon_r} - 1\right)
\end{equation}
\begin{equation}
  \ln\hat{\phi}_i^\mathrm{Born} = -\frac{e^2 z_i^2}{8\pi\varepsilon_0 k_B T R_{B,i}}
    \left(\frac{1}{\varepsilon_r} - 1\right)
    + \frac{e^2}{8\pi\varepsilon_0 k_B T}\sum_j \frac{n_j z_j^2}{R_{B,j}}
    \frac{1}{\varepsilon_r^2}\frac{\partial \varepsilon_r}{\partial n_i}
\end{equation}

\noindent The relative permittivity $\varepsilon_r$ is computed from the water fraction and temperature
via the Johnson--Uematsu--Franck correlation, with a Clausius--Mossotti correction for the CO$_2$
contribution using its dipole polarizability.

\end{document}
